\scanpage{27}
% archive sheet number: 26
\mspage{114}
% The following draft is upside down in the scan and cancelled by a long diagonal stroke.
 \deleted{On considère alors le groupe
\[
 \widetilde Z=G/G_0^+\qquad\text{avec}\qquad G\xrightarrow{\tilde\delta}\widetilde Z \tag*{(9)}
\]
qui est un rev. de degré 2 de $Z$
\[
 1\longrightarrow\{\pm1\}\xrightarrow[\unsure{\prime|,}]{}  \widetilde Z\longrightarrow Z\longrightarrow 1,\tag*{(10)}
\]
et le diagramme d'inclusions de groupes
\diagram{\begin{tikzcd}[ampersand replacement=\&,column sep=small,row sep=small]
 {} \& U\arrow[r,hook,no head,"Z","{\deleted{\unsure{\mathrm{dist.}}}}"']\arrow[d,hook,no head,"\mu_2"'] \& B''\arrow[r,hook]\arrow[d,hook,no head,"\mu_2"'] \& G_0^+\arrow[d,hook,no head,"\mu_2"'] \& {}\\
\widetilde Z\arrow[u]\arrow[d] \& \widetilde U\arrow[r,hook,no head,"Z","{\deleted{\unsure{\mathrm{dist.}}}}"']\arrow[d,hook,"Z"'] \& \widetilde B''\arrow[r,hook]\arrow[d,hook,"Z"'] \& G_0\arrow[d,hook,"Z"'] \& \widetilde Z\arrow[u]\arrow[d]\\
{} \& B'\arrow[r,hook,no head,"Z","{\deleted{\unsure{\mathrm{dist.}}}}"'] \& B\arrow[r,hook,no head] \& G \& {}
\end{tikzcd}\quad (11)}
 \margin{On a écrit au dessus de toute flèche d'inclusion d'un sous-groupe dist. dans un groupe \unsure{ambiant}, le groupe quotient correspondant (à isom. can. près)}
 \deleted{\unsure{$\bigg\}$}}, 
défini ainsi à partir de
\[
 G_0^{+}\subset G_0\subset G
\]
et de l'action de $G$ sur $S$, en termes d'un $G_0^+$-répére
\[
 (s_0,t_0)\in\overrightarrow{A_0}\qquad\text{i.e.}\quad s_0,t_0\in S,\quad\{s_0,t_0\}\in A_0:\tag*{(12)}
\]
\[
\left\{\begin{aligned}
 B'&=G_{s_0}=\{g\in G\mid gs_0=s_0\}\\
 \widetilde U&=B'\cap G_0\deleted{\unsure{{}_{s_0}}}=\{g\in G_0\mid gs_0=s_0\}=\\
 &=\{b'\in B'\mid\delta(b')=1\}\\
 U&=B'\cap G_0^+=\{g\in G_0^+\mid gs_0=s_0\}=\{b'\in B'\mid\tilde\delta(b')=1\}\\
 \widetilde B''&=\{g\in G_0\mid gs_0\in Z.s_0\}\leftarrow\\
 B''&=\widetilde B''\cap G_0^+\quad\deleted{\unsure{\{}}
 \quad[\text{NB }\widetilde B''=\widetilde U.B'']\\
 B&=B'.B''\,(=B'.\widetilde B'')\quad\deleted{\text{NB }\unsure{\{}\,\widetilde B''=\widetilde U.B''}.
\end{aligned}\right.\tag*{(13)}
\]
 \margin{c'est le sous-groupe de $\operatorname{Norm}_{G_0}(\widetilde U)$ contenant $\widetilde U=G_{0s_0}$ qui sert à décrire $Z\simeq\widetilde B''/\widetilde U$ et son action sur $S\simeq G_0/\widetilde U$}
 \margin{\uline{NB} $B$ est bien un \unsure{ss-groupe|sous-groupe|s-groupe} car $B'=G_{s_0}$ est \unsure{stabilisé|normalisé|normalise} par $B''$ qui est défini en termes de $s_0$}
 \margin{produits des sous-ens. dans un groupe $G$\quad $A.B=\{ab\mid a\in A,b=B\}$}}
